\newcommand{\graph}[1][k]{G_{#1}}

In this section, we prove that the bound in \cref{lemma:twin-cover-hcupn} is tight. To show this, for every $k \in \N$ with $k \geq 2$, we build a graph $\graph$ of $2^k + k$ vertices such that $\clustermc(\graph) = 2^{\vc(\graph)} + \vc(\graph)$. As every vertex cover is also a twin-cover, this will prove the tightness of \cref{lemma:twin-cover-hcupn}.

The graph $\graph$ is constructed as follows.
Let $X$ be a set of $k$ vertices, and for each subset $X' \subseteq X$, let $v_{X'}$ be a vertex. We define
\begin{align*}
    V(\graph) &\coloneq X \cup \{v_{X'} \mid X' \subseteq X\}, \text{ and}\\
    E(\graph) &\coloneq \{(x, v_{X'}) \mid x \in X'\}.
\end{align*}

In other words, we create a set $X$ of vertices, and create one extra vertex $v_{X'}$ for each possible neighborhood $X'$ inside $X$. Two vertices outside $X$ are never adjacent to each other, and neither are two vertices inside $X$. Observe that this graph has $2^k + k$ vertices, as there are $2^k$ possible subsets $X' \subseteq X$. \Cref{fig:G-k} shows the graph $\graph[3]$.

\begin{figure}[ht]
    \centering
\begin{tikzpicture}[xscale=1.8,yscale=1.1]
    \draw[gray,fill=gray!10,rounded corners=5] (2.4,2.7) rectangle (8.6,3.75);
    \node[gray!40!black] at (2.15,3.225) {$X$};

    \node[style=nodo,label=above:{$x_1$}] (x1) at (3.5,3) {};
    \node[style=nodo,label=above:{$x_2$}] (x2) at (5.5,3) {};
    \node[style=nodo,label=above:{$x_3$}] (x3) at (7.5,3) {};

    \node[style=nodo,label=below:{$v_{\emptyset}$}] (v0) at (2,0) {};
    \node[style=nodo,label=below:{$v_{\{x_1\}}$}] (v1) at (3,0) {};
    \node[style=nodo,label=below:{$v_{\{x_2\}}$}] (v2) at (4,0) {};
    \node[style=nodo,label=below:{$v_{\{x_3\}}$}] (v3) at (5,0) {};
    \node[style=nodo,label=below:{$v_{\{x_1,x_2\}}$}] (v12) at (6,0) {};
    \node[style=nodo,label=below:{$v_{\{x_1,x_3\}}$}] (v13) at (7,0) {};
    \node[style=nodo,label=below:{$v_{\{x_2,x_3\}}$}] (v23) at (8,0) {};
    \node[style=nodo,label=below:{$v_{\{x_1,x_2,x_3\}}$}] (v123) at (9,0) {};

    \foreach \v in {v1,v12,v13,v123} \draw[style=arista] (x1) -- (\v);
    \foreach \v in {v2,v12,v23,v123} \draw[style=arista] (x2) -- (\v);
    \foreach \v in {v3,v13,v23,v123} \draw[style=arista] (x3) -- (\v);
\end{tikzpicture}
    \caption{The graph $\graph[3]$.}
    \label{fig:G-k}
\end{figure}

The main property of the graph $\graph$ is that no two vertices have the same neighborhood, either closed or open.
\begin{lemma}
    \label{lemma:no-twins-in-G}
    For every edge $(x, v_{X'}) \in E(\graph)$, we have that $N[x] \neq N[v_{X'}]$.
\end{lemma}
\begin{proof}
    As $k \geq 2$, each vertex $x \in X$ has at least two neighbors outside of $X$. On the other hand, $v_{X'}$ has no neighbors outside of $X$. Hence $x$ has a neighbor that $v_{X'}$ does not have.
\end{proof}

\begin{lemma}
    \label{lemma:no-false-twins-in-G}
    For every pair of non-adjacent vertices $u, w \in V(\graph)$, we have that $N(u) \neq N(w)$.
\end{lemma}
\begin{proof}
    If $u, w \notin X$, they have different neighborhoods inside $X$, and thus $N(u) \neq N(w)$. On the other hand, if one of them, say $u$, belongs to $X$, then the vertex $v_{\{u\}}$ is adjacent to $u$ and not to $w$.
\end{proof}

With these two lemmas, we close this section by determining $\vc(\graph)$ and $\clustermc(\graph)$.
\begin{lemma}
    \label{lemma:vc-G}
    $\vc(\graph) = k$.
\end{lemma}
\begin{proof}
    Note that every edge in $\graph$ is incident to a vertex in $X$, making $X$ a vertex cover of $\graph$. Hence, $\vc(\graph) \leq k$.

    By the definition of vertex cover, every edge in $\graph$ must have an endpoint in every vertex cover of $\graph$. In particular, the matching of size $k$ formed by the edges $(x, v_{\{x\}})$ for each $x \in X$ shows that every vertex cover must have size at least $k$. Thus, we also have $\vc(\graph) \geq k$, which gives the statement of the lemma.
\end{proof}

\begin{lemma}
    \label{lemma:clustermc-G}
    $\clustermc(\graph) = 2^k + k$.
\end{lemma}
\begin{proof}
    As $|V(\graph)| = 2^k + k$, this is equivalent to proving that the only $\cluster$-modular partition of $\graph$ is the one that has one part per vertex.

    Suppose, to the contrary, that there exists a $\cluster$-modular partition $\mathcal{M}$ of $\graph$ into fewer parts. Then some module $M \in \mathcal{M}$ contains two distinct vertices $u$ and $w$, which are either adjacent or not. We show that both cases lead to a contradiction.

    Suppose first that $u$ and $w$ are adjacent, so they are the two endpoints of an edge $(x, v_{X'})$. By \cref{lemma:no-twins-in-G}, these two vertices are not twins. As every pair of adjacent vertices in a part of a $\cluster$-modular partition must be twins, this yields a contradiction.

    Otherwise, $u$ and $w$ are not neighbors. By \cref{lemma:no-false-twins-in-G}, one of them, say $u$, has a neighbor $z$ that is not a neighbor of $w$. By the previous paragraph, every pair of adjacent vertices belong to different parts. Hence, $z \notin M$. Therefore, the two vertices $u$ and $w$ in $M$ have different neighborhoods outside $M$, and thus $M$ is not a module. This contradicts the definition of $\cluster$-modular partition.
\end{proof}

Combining \cref{lemma:vc-G,lemma:clustermc-G}, we obtain that $\clustermc(\graph) = 2^{\vc(\graph)} + \vc(\graph)$ for every $k \geq 2$, and thus the bound in \cref{lemma:twin-cover-hcupn} is tight, even when replacing the twin-cover number with the vertex cover number.